% generated by roam.py handbook export -- не редактировать, править заметки
% 20260920221510-критерии_выпуклости_функции_однои_переменнои_через_вторую_производную.org (body 02f42c324a24)

\begin{theorem}[{Критерий выпуклости функции одной переменной через вторую производную}]\label[theorem]{note:9CEE854D-ED37-4C08-82EB-E41B36C43EA6}
  Пусть $I \subset \R$~--- промежуток, $\phi: I \rightarrow \R$~--- дважды дифференцируемая функция.
  Тогда $\phi$~--- \hyperref[note:1CA14F2C-637F-41DE-9E97-72075D19E97F]{выпуклая} $\Leftrightarrow$ $\phi^{\prime\prime}(t) \ge 0\ \forall t \in I$
\end{theorem}

\begin{proof}
  Оба направления идут через промежуточное условие: $\phi^{\prime}$ не убывает на $I$. Для дифференцируемой $\phi^{\prime}$ это равносильно $\phi^{\prime\prime} \ge 0$ (признак монотонности). Поэтому достаточно доказать: $\phi$~--- выпуклая $\Leftrightarrow$ $\phi^{\prime}$ не убывает.

  \proofpart{$\Rightarrow$}
  Пусть $\phi$~--- выпуклая, $s, t \in I$. Запишем неравенство из \hyperref[note:2F475EE4-2630-45FD-8DD2-8EA9D48059E4]{одномерного дифференциального критерия} дважды:
  \[
  \phi(t) \ge \phi(s) + \phi^{\prime}(s)(t-s)
  \]
  \[
  \phi(s) \ge \phi(t) + \phi^{\prime}(t)(s-t)
  \]
  Сложим их и сократим $\phi(s) + \phi(t)$:
  \[
  0 \ge \phi^{\prime}(s)(t-s) + \phi^{\prime}(t)(s-t) = -(\phi^{\prime}(t) - \phi^{\prime}(s))(t-s)
  \]
  То есть $(\phi^{\prime}(t) - \phi^{\prime}(s))(t-s) \ge 0$: при $t > s$ получаем $\phi^{\prime}(t) \ge \phi^{\prime}(s)$~--- производная не убывает.

  \proofpart{$\Leftarrow$}
  Пусть $\phi^{\prime}$ не убывает, $s, t \in I$, $s \ne t$. По теореме Лагранжа существует $\xi$ между $s$ и $t$:
  \[
  \phi(t) - \phi(s) = \phi^{\prime}(\xi)(t-s)
  \]
  Если $t > s$, то $\xi > s \Rightarrow \phi^{\prime}(\xi) \ge \phi^{\prime}(s)$, и, умножая на $t - s > 0$, получаем $\phi^{\prime}(\xi)(t-s) \ge \phi^{\prime}(s)(t-s)$.
  Если $t < s$, то $\xi < s \Rightarrow \phi^{\prime}(\xi) \le \phi^{\prime}(s)$, и, умножая на $t - s < 0$, снова получаем $\phi^{\prime}(\xi)(t-s) \ge \phi^{\prime}(s)(t-s)$.
  В обоих случаях
  \[
  \phi(t) \ge \phi(s) + \phi^{\prime}(s)(t-s)
  \]
  а это по \hyperref[note:2F475EE4-2630-45FD-8DD2-8EA9D48059E4]{одномерному дифференциальному критерию} означает выпуклость $\phi$.
\end{proof}
